\part{Aktueller Stand}
\section{Basisarchitektur und -funktionalitäten des cores}

\begin{frame}{Basisarchitektur und -funktionalitäten des \texttt{core}s}{Aktueller Stand}
  \begin{itemize}
    \item Support für alle drei wesentlichen IO-Multiplexer
    \begin{itemize}
      \item \texttt{select()}
      \item \texttt{poll()}
      \item \texttt{epoll()}
    \end{itemize}
    \item \texttt{SingleShot} und \texttt{IntervalTimer} Klassen
    \item Dateizugriff über Klasse \texttt{File} (zur Zeit nur lesen)
    \item Template-Basisklassen für konkrete \texttt{SocketServer} und \texttt{SocketClient} Klassen
    \item Abstrakte Basisklasse \texttt{SocketContext} für Applikationsprotokolle
    \item Transparente Unterstützung von \texttt{ssl/tls} für \texttt{SocketServer} und \texttt{SocketClient} Templates
    \item Dynamisches Laden von Plugins (\texttt{dlopen()}, \texttt{dlsym()}, \texttt{dlclose()})
    \begin{description}[WS-Subprotocols:]
      \item[HTTP-Upgrade:] \texttt{websocket} (etwas Anderes ist (noch) nicht spezifiziert)
      \item[WS-Subprotocol:] \texttt{tiktaktoe}, \texttt{echo}, \texttt{mqtt}
    \end{description}
    \item Daemonize wenn gewünscht
    \item Automatisches Bereitstellen von Kommandozeilenparameter und
    \item Unterstützung von Konfigurationsdateien
  \end{itemize}
\end{frame}
